\section{Verificación del tablero operativo}\label{sec:pruebas-tablero}

La aceptación del agregado separa autorización, consistencia y presentación. Sus comprobaciones deben distinguir el ámbito completo de Owner del conjunto asignado a Litigante, rechazar Paralegal y Cliente, y evitar que una pertenencia retirada conserve totales o carga de trabajo. La fecha común de consulta y la ausencia de almacenamiento en caché permiten identificar la observación, sin convertirla en un cálculo jurídico nuevo.

\rutarepo{scripts/api-dashboard-demo.py}, integrado en la demostración HTTP, captura los contadores y la carga de trabajo después de preparar los datos de prueba. La restauración de PostgreSQL y el nuevo acceso con MFA se realizan mediante el procedimiento existente; a continuación se repite la consulta y se compara todo el agregado salvo el instante de observación. Ambos instantes se validan como lecturas UTC nuevas. Los plazos sintéticos del recorrido utilizan fechas de enero y febrero de 2026; un cambio de contador, incluso por cruzar una frontera temporal al modificar esos datos, produce un fallo explícito y no se descarta de la comparación. Las denegaciones y el ámbito vacío del litigante revocado también se vuelven a comprobar.

La aceptación integrada aprobó el 27 de septiembre de 2026: los contadores y la carga restaurados fueron idénticos, con nueva lectura UTC, MFA renovado y denegaciones conservadas. Las comprobaciones focales aprobaron cinco casos de aplicación, ocho de PostgreSQL y tres de HTTP; la interfaz aprobó siete verificaciones de cliente, ocho recorridos controlados y uno con servicios reales que retiró una asignación y comprobó el agregado vacío. Se inspeccionaron las vistas de escritorio y móvil. Los resultados nuevos se registran en \rutarepo{docs/verification-report.md}; las campañas históricas no acreditan por sí mismas el comportamiento añadido.

\section{Verificación de la creación contextual de plazos}\label{sec:pruebas-plazo-contextual}

La aceptación de este flujo debe comprobar que preparar no crea ninguno de los dos registros y que confirmar conserva una única operación, sus recibos y las capturas exactas seleccionadas. Incluye repetición idéntica, denegación por cambio de permisos o dependencias, rechazo de operaciones ordinarias preexistentes y concurrencia. La atomicidad exige además provocar un fallo durante la transacción y comprobar que no quedan escrituras parciales; rechazar una precondición antes de escribir no demuestra por sí solo esa propiedad.

El 27 de septiembre de 2026 la demostración HTTP integrada terminó con salida cero: comprobó confirmación conjunta, repetición exacta, consulta en agenda y rechazo de una preparación obsoleta. Después de restaurar PostgreSQL y renovar la sesión, conservó las capturas, los recibos del plazo y la asociación, y su proyección en agenda. Las comprobaciones focales del backend aprobaron cinco casos de aplicación, nueve distintos de PostgreSQL y tres de HTTP; no constituyen una regresión global. La interfaz aprobó once verificaciones Node, once recorridos con HTTP controlado en 28.8 s y uno con servicios reales en 13.6 s de escenario (17.8 s del ejecutor de navegador). Se inspeccionaron las capturas de escritorio y móvil, sin desbordamiento horizontal. La regresión global permanece pendiente. La compilación y revisión del PDF se registran por corte en el informe de verificación. La evidencia de la \cref{sec:impl-plazo-contextual} se registra por separado del tablero en \rutarepo{docs/verification-report.md}.
